\cotefile{12}
\scanpage{137}
\[
R^{\deleted{\unsure{*}}\unsure{*|n|\bullet}}(if)_* \Leftarrow R^p i_*(R^q f_*(\mathbb Q_\ell))
\]
\diagram{\begin{tikzcd}[ampersand replacement=\&,execute at end picture={
\draw (-.42,1.05)--(3.86,1.05)--(3.92,.21)--(-.43,.29);
\draw (0,1.32)--(.04,.29);
\draw (.02,.65)--(3.85,.65);
\node[left] at (-.22,1.07) {$2$};
\node[below right] at (3.76,.21) {$N$};
\fill (1.37,1.05) circle (1.5pt);
\foreach \x in {1.36,1.52,1.60,1.96,2.29,2.59}{\fill (\x,.65) circle (1.6pt);}
\fill (2.60,.24) circle (1.7pt);
\draw[->,thick] (1.40,1.03)--(1.91,.74);
\draw[->,thick] (2.04,.57)--(2.57,.31);
}]\phantom{.}\end{tikzcd}}
avec $E^{pq}_2=0$ si $q\notin[0,2]$

\unsure{D'ailleurs|D'ailleur} $\unsure{\mathbb R|R|\mathbf R}^0 f_*(\mathbb Q_\ell)$ est essentiellement $\mathbb Q_\ell$ (au moins si $f$ à fibres connexes, auquel cas on se ramène facilement), donc les \deleted{$R^p i_*$ sont} $E^{p0}_2$ sont $(d-1)$-admissibles par \deleted{l'hyp.} ce qu'on a déjà \unsure{prouvé|vu}.

\deleted{Donc} D'autre part $R^2 f_*(\mathbb Q_\ell)\simeq\mathbb Q_\ell(-1)$, et
\[
 E^{p2}_2\simeq R^p i_*(R^2 f_*(\mathbb Q_\ell))\simeq
 \deleted{\unsure{E^{p2}_2}\ \unsure{R^p i_*}}\inserted{\mathbb Q_\ell(-1)}\otimes E^{p0}_2
\]
est donc également $(d-1)$ admissible, à fortiori \deleted{$\unsure{d+d'-1|d+d'-0}$}\inserted{$(d+d'-1)$-}admissible\inserted{ i.e. $d$-admissible}. \deleted{Pour} On en conclut, utilisant lemme 2\inserted{ et la constructibilité des $E^{pq}_2$}, que $E^{p1}_2$ \deleted{\unsure{$=$}} est une extension de \deleted{\unsure{deux}}\inserted{trois} faisceaux, les extrêmes deux \deleted{\unsure{$(d+d'-1)$-admis}} $d$-admissibles, et le médian étant $E^{p1}_3=\ldots=E^{p1}_\infty$. Il suffit donc de prouver que $E^{p1}_\infty$ est $d$-admissible, ou encore que $R^{\unsure{n|p1|p+1}}(if)_*(\mathbb Q_\ell)$ l'est, ou encore que $R^n(gf)_*(\unsure{X'|\chi'|X''})$ l'est, ce qui par Leray et lemme 3 nous ramène au cas des $R^{\unsure{-|n|\cdot}}j_*(\mathbb Q_\ell)$, qu'on sait traiter\ldots
% archive sheet number: 136
